
\documentclass[aspectratio=169]{beamer}

\mode<presentation>
\usetheme{Boadilla}
\definecolor{NTHU1}{RGB}{127,16,132}
\definecolor{NTHU2}{RGB}{228,210,231}
\definecolor{blue}{RGB}{30,90,205}
\definecolor{red}{RGB}{213,94,0}
\definecolor{green}{RGB}{0,128,0}
\definecolor{charcoalgray}{RGB}{108,104,104}
\setbeamercolor{title}{fg=NTHU1}
\setbeamercolor{frametitle}{fg=NTHU1}
\setbeamercolor{block title}{bg=NTHU1, fg=white}
\setbeamercolor{block body}{bg=NTHU2}
\setbeamercolor{structure}{fg=NTHU1}
\setbeamercolor{item projected}{fg=white}
\setbeamercolor{item}{fg=NTHU1}
\setbeamercolor{subitem}{fg=NTHU1}
\setbeamercolor{section in toc}{fg=NTHU1}
\setbeamercolor{description item}{fg=NTHU1}
\setbeamercolor{caption name}{fg=NTHU1}
\setbeamercolor{button}{bg=NTHU1, fg=white}
\setbeamercolor{caption name}{fg=NTHU1}
\usepackage{graphics}
\usepackage{tikz}
\usepackage{amsmath}
\usepackage{bbm}
\usetikzlibrary{decorations.pathreplacing}
\usepackage{geometry}
\usepackage{booktabs}
\usepackage{bookmark}
\usepackage{multirow, makecell}
\usepackage{float}
\usepackage{fancyvrb}
\usepackage{caption}
\captionsetup[figure]{singlelinecheck = false, format= hang, justification=centering}
\captionsetup[subfigure]{singlelinecheck = false, format= hang, justification=raggedright}
\usepackage{subcaption}
\usepackage{adjustbox}
\usepackage{hyperref}
\usepackage{threeparttable}
\usepackage{appendixnumberbeamer} 
\usepackage[T1]{fontenc}
\usepackage[default]{lato} 
\usepackage{smartdiagram}
\smartdiagramset{
  back arrow disabled = true,
  module minimum width=1cm,
  module x sep = 3,
  uniform arrow color=true,
}
\usepackage{xcolor}
\usepackage{colortbl}
	\definecolor{light-gray}{gray}{0.99}

\newenvironment{wideitemize}{\itemize\addtolength{\itemsep}{10pt}}{\enditemize}
\newenvironment{wideenumerate}{\enumerate\addtolength{\itemsep}{10pt}}{\endenumerate}
\newenvironment{widedescription}{\description\addtolength{\itemsep}{10pt}}{\enddescription}
\hypersetup{
colorlinks=true,
linkcolor=NTHU1,
filecolor=green, 
urlcolor=blue,
}
\beamertemplatenavigationsymbolsempty
\setbeamercolor{author in head/foot}{bg=white, fg=NTHU1}
\setbeamercolor{title in head/foot}{bg=white, fg=NTHU1}
\setbeamercolor{date in head/foot}{bg=white, fg=NTHU1}
\setbeamercolor{section in head/foot}{bg=white, fg=NTHU1}
\setbeamercolor{headline}{bg=white}
\setbeamertemplate{footline}{
    \leavevmode%
    \hbox{%
        \begin{beamercolorbox}[wd=.333333\paperwidth,ht=2.25ex,dp=1ex,center]{date in head/foot}%
            \usebeamerfont{date in head/foot}%\insertdate
        \end{beamercolorbox}%
        \begin{beamercolorbox}[wd=.333333\paperwidth,ht=2.25ex,dp=1ex,center]{title in head/foot}%
            \usebeamerfont{title in head/foot}\insertshorttitle
        \end{beamercolorbox}%
        \begin{beamercolorbox}[wd=.333333\paperwidth,ht=2.25ex,dp=1ex,right]{date in head/foot}%
            \usebeamerfont{date in head/foot}\insertframenumber{}\hspace*{2em}
        \end{beamercolorbox}}%
        \vskip0pt%
    }
    %% May need to script the introduction!!!!!!
%\setbeamercolor{page number in head/foot}{fg=NTHU1}
\setbeamertemplate{section in toc}[sections numbered]
\setbeamertemplate{subsection in toc}{\leavevmode\leftskip=3em\rlap{\hskip-1.75em\inserttocsectionnumber.\inserttocsubsectionnumber}
\inserttocsubsection\par}
\setbeamerfont{subsection in toc}{size=\footnotesize}
%\setbeamertemplate{headline}{%
  %\begin{beamercolorbox}[ht=5.5ex]{section in head/foot}
    %\vskip2pt\insertnavigation{0.33\paperwidth}\vskip2pt
  %\end{beamercolorbox}%
%}
\newenvironment{transitionframe}{\setbeamercolor{background canvas}{bg=white}\setbeamertemplate{footline}{} \begin{frame}[noframenumbering]}{\end{frame}}


\makeatletter
\let\@@magyar@captionfix\relax
\makeatother


\title[]{Public Economics and Development} % Change this regularly
\subtitle[]{Week 13: Recruiting Capable and Motivated Individual to Public Sector}
\author[Lee]{Seung-hun Lee }
\institute[]{Taipei School of Economics\\ National Tsing Hua University}

\date[]{May 19th, 2026}

\begin{document}
\begin{transitionframe}
\titlepage
\end{transitionframe}


%%% Color slides for section headers: Use for colloquium version (The ones bewteen \iffals and \fi)

\section{Overview setup of the topic}
\begin{frame}
\frametitle{Why Selection Matters: The Stakes}
Patronage and meritocracy coexist in most countries
\begin{itemize}\setlength\itemsep{3pt}
\item Low-skilled positions: patronage common (reward political supporters)
\item High-skilled positions: meritocracy more likely (incompetence could be costly)
\end{itemize}\medskip
Quality of governance strongly linked to \textbf{how} bureaucrats and personnel are selected
\begin{itemize}\setlength\itemsep{3pt}
\item But selection is often distorted — patronage and political connections pervasive
\item Cost: Loyalty over capability — disrupts continuity and undermines service delivery
\begin{itemize}\setlength\itemsep{3pt}
\item Junior bureaucrats under-invest in skills when transfer risk is high
\item Inconsistent and inefficient practices within public organizations
\end{itemize}
\item Benefit: Discretion can weed out entrenched, low-performing workers
\begin{itemize}\setlength\itemsep{3pt}
\item May be necessary when civil service protections shield underperformers
\item Allows principal to restore control over misaligned bureaucrats 
\end{itemize}
\end{itemize}
\end{frame}


\begin{frame}
  \frametitle{Stylized facts on political influence in recruitment}
  \begin{figure}
    \centering
    \includegraphics[keepaspectratio, width=0.85\textwidth]{fig1.png}
  \end{figure}
\end{frame}

\begin{frame}
\frametitle{What Determines Selection Quality?}
Rule-based vs. discretionary selection — a fundamental tradeoff
\begin{itemize}\setlength\itemsep{3pt}
\item \textbf{Merit-based}: reduces favoritism, improves average quality, builds institutional trust
\begin{itemize}\setlength\itemsep{3pt}
\item Civil service exams, competitive hiring, transparent criteria
\item But may not select for motivation, integrity, or adaptability
\end{itemize}
\item \textbf{Discretion-based}: allows flexibility in matching candidates to roles
\begin{itemize}\setlength\itemsep{3pt}
\item Can select for unobservable traits — motivation, integrity, cultural fit
\item But opens door to favoritism, patronage, and political interference
\end{itemize}
\end{itemize}\medskip
Financial incentives shape who applies — but not always in the right direction
\begin{itemize}\setlength\itemsep{3pt}
\item Higher wages attract more and better candidates
\item May signal wrong job characteristics and crowd out intrinsically motivated candidates
\end{itemize}\medskip
Manager selection is often the missing link
\begin{itemize}\setlength\itemsep{3pt}
\item Managers determine practices within organizations
\item Yet managerial positions often exempt from meritocratic reforms
\end{itemize}
\end{frame}







\begin{transitionframe}
\frametitle{Roadmap for today}
\tableofcontents
\end{transitionframe}








\section{Impact of management on the efficiency of public organizations}


\subsection{Dal Bó, Finan, Folke, Persson, Rickne (2017): Who Becomes a Politician?}
\begin{transitionframe}
\begin{center}
  \begin{huge}
    \textcolor{NTHU1}{%
Dal Bó, Finan, Folke, Persson, Rickne (2017):\\ Who Becomes a Politician?
    }
  \end{huge}
\end{center}
\end{transitionframe}

\begin{frame}
\frametitle{Q: How to attract competent leaders while being representative?} 
Can democracies form competent and representative leadership (inclusive meritocracy)?
\begin{itemize}\setlength\itemsep{3pt}
\item Voters want to elect leaders capable of efficiently implementing policies
\item Voters also want policy-makers who represent diverse interests
\item Classical economic models suggest that these two are in a trade-off relation
\begin{itemize}\setlength\itemsep{3pt}
\item Market wage < Returns from office for low-quality individuals (low opportunity cost)
\item Good politicians likely to `free-ride' on bad politicians $\to$ Supply of good politicians $\Downarrow$
\item Even if good candidates can be selected, may make broad representation difficult
\end{itemize}
\end{itemize}\medskip
This paper uses data on Swedish local politicians to identify patterns in inclusive democracy
\begin{itemize}\setlength\itemsep{3pt}
\item Includes elected and \textit{nonelected}: Address selection bias 
\item Detailed measures on quality and representation
\item Identifies how both competent and representative leaders are selected in democracies
\end{itemize}
\end{frame}

\begin{frame}
\frametitle{Context and data} 
Swedish National MPs and Municipal council members
\begin{itemize}\setlength\itemsep{3pt}
\item MPs are full-time, while most municipal council members have other work
\begin{itemize}
\item Mayors, selected among council members, are full-time and paid highly
\end{itemize}
  \item Low monetary cost to individuals running for local office
  \item Key motives: Intrinsic concerns with policy, social interaction in policy circles
  \item Left represents blue-collar, center-right represents white-collar workers
\end{itemize}\medskip
Data obtained from 
\begin{itemize}\setlength\itemsep{3pt}
\item  Sample: All candidates from 1980-2010 
\item Links candidates to admin registers (including military enlistment and tax records)
\item Age, sex, education level, earnings, occupation, parental/sibling info, cognitive capacity 
\end{itemize}
\end{frame}

\begin{frame}
\frametitle{Positive selection of politicians based on abilities} 
\begin{figure}
  \includegraphics[keepaspectratio, width=0.65\textwidth]{dffpr2017_f1.png}
\end{figure}
\begin{itemize}
\item Beats classical prediction that there are more low-quality individuals entering politics
\end{itemize}
\end{frame}


\begin{frame}
\frametitle{Are they socially representative?}
Determining source of positive selection
\begin{itemize}\setlength\itemsep{3pt}
\item Privileged access? Ability should matter little once elite status is controlled for %$ By the heirs of rich families
\item Exclusive meritocracy? Key abilities are strongly associated with social background
\item Inclusive meritocracy? High ability individuals come from all social backgrounds
\end{itemize}\medskip
Empirical tests
\begin{itemize}\setlength\itemsep{3pt}
\item Compare with siblings: Same "elitist background" $\to$ Positive selection still appears
\item Parental income: Social background of politicians is evenly distributed
\item Social class: Politicians from wide range of occupations represented
\end{itemize}
\end{frame}

\begin{frame}
\frametitle{Elitism unlikely to explain positive selection} 
\begin{figure}
  \includegraphics[keepaspectratio, width=0.7\textwidth]{dffpr2017_f2.png}
\end{figure}
\end{frame}

\begin{frame}
\frametitle{Neither does exclusive meritocracy} 
\begin{figure}
  \includegraphics[keepaspectratio, width=0.8\textwidth]{dffpr2017_f3.png}
\end{figure}
\end{frame}

\begin{frame}
\frametitle{Discussion and takeaways}
What drives competence and representation? 
\begin{itemize}\setlength\itemsep{3pt}
\item Weak trade-off between competence and social background
\item Positive selection on ability even among those with poorer backgrounds
\item Intrinsic motives: Positive selection even for low-paid municipal councils
\item Screening by parties: Sorts more able people to higher positions
\end{itemize}\medskip
Remaining questions
\begin{itemize}\setlength\itemsep{3pt}
\item When and where do these findings hold outside 1980-2010 Sweden?
\begin{itemize}\setlength\itemsep{3pt}
\item Intrinsically motived individuals may not enter politics elsewhere, how to attract them?
\item Party incentives: Promote someone who can win vs. who has best policy
\item Other democratic countries? Sweden now? 
\end{itemize}
\end{itemize}
\end{frame}


\subsection{Akhtari, Moreira, Trucco (2022): Political Turnover, Bureaucratic Turnover, and the Quality of Public Service}
\begin{transitionframe}
\begin{center}
  \begin{huge}
    \textcolor{NTHU1}{%
Akhtari, Moreira, Trucco (2022):\\ Political Turnover, Bureaucratic Turnover, and the Quality of Public Service
    }
  \end{huge}
\end{center}
\end{transitionframe}




\begin{frame}
\frametitle{Q: How do changes in political leadership affect public service?} 
Countries vary in the extent to which politicians control appointment of personnel
\begin{itemize}\setlength\itemsep{3pt}
\item East Asia: Only the top chiefs and vice chiefs are politically appointed
\item On the other end, much more positions are appointed by the president
  \begin{itemize}\setlength\itemsep{3pt}
  \item Nearly all of top 500 government agency positions appointed in Brazil, Israel, Nigeria etc.
  \item 26\% of all govt personnel work under some threat of removal by the president
\end{itemize}
\item Thus, politicaly motivated replacements are commonplace in the latter 
\end{itemize}\medskip
This paper examines the effects of such scope of politically motivated appointments
\begin{itemize}\setlength\itemsep{3pt}
  \item Leverage bureaucratic turnover following close elections in public schools
  \item Local politicians have significant influence over public school administration
\item Turnovers in political leadership $\to$ New personnel $\to$ Education peformance $\downarrow$
\item Documents trade-off between capability and loyalty 
\end{itemize}
\end{frame}



\begin{frame}
\frametitle{Municipality-run schools in Brazil} 
Municipal education in Brazil
\begin{itemize}\setlength\itemsep{3pt}
\item 5,565 decentralized municipalities are responsible for essential services
\item Mayors are elected every 4 years, and can appoint $\sim$30\% of workers via `contract'
\begin{itemize}\setlength\itemsep{3pt}
\item Rest have passed civil service exams
\end{itemize}
\item Municipalities responsible for daily operations and personnel of primary schools
\begin{itemize}\setlength\itemsep{3pt}
\item 60\% of headmasters in public schools are politically appointed
\item 35\% of teachers hired by contract (rest passed civil service exams)
\end{itemize}
\end{itemize}\medskip
Why local politicians care about having loyalists in managerial positions
\begin{itemize}\setlength\itemsep{3pt}
\item Managers can use control over public programs and resources politically
\item Headmaster position used as means to provide contracts to political supporters
\item Thus, public positions used as means to consolidate political support
\end{itemize}
\end{frame}



\begin{frame}
  \frametitle{Data and empirical strategies}
Data sources: Brazilian election and education data (`08-`12)
\begin{itemize}\setlength\itemsep{3pt}
\item Municipal personnel changes from Registry of Social Information (RAIS)
\item National exam score (Prova Brasil) and School Census for edu quality and personnel
\item Electoral data (party, vote share) from Tribunal Superior Eleitoral (TSE)
\item Municipal characteristics from various geographical sources and political atlas
\end{itemize}\medskip
Using municipalities with close election results in regression discontinuity
\begin{itemize}\setlength\itemsep{3pt}
\item Not just comparing places with turnovers vs without one
\begin{itemize}\setlength\itemsep{3pt}
\item Incompetent incumbent (bad public goods) $\to$ Voted out $\to$ reverse causality
\end{itemize}
\item Narrow down to places with and without turnovers AND small electoral margins
\[
Y_{imt+1}=\alpha + \beta I[\text{Turnover}]_{mt}+ \gamma \text{Margins}_{mt} + \delta  I[\text{Turnover}]_{mt}\times\text{Margins}_{mt} + \lambda X_{jmt} + e_{jmt}
\]\vspace{-20pt}
\item Municipalities included in RD similar except for turnover, and no sorting on margins
\end{itemize}
\end{frame}

\begin{frame}
  \frametitle{Political turnover $\to$ Replacement of municipal personnel}
  \begin{figure}
    \centering
  \includegraphics[keepaspectratio, width=0.7\textwidth]{amt2022_f1.png}
  \end{figure}
\end{frame}



\begin{frame}
  \frametitle{Lower educational performance of students}
  \begin{figure}
    \centering
  \includegraphics[keepaspectratio, width=1\textwidth]{amt2022_f2.png}
  \end{figure}
\end{frame}

\begin{frame}
  \frametitle{Drastic changes in school personnel (managerial and frontline)}
  \begin{figure}
    \centering
  \includegraphics[keepaspectratio, width=1\textwidth]{amt2022_f3.png}
  \end{figure}
\end{frame}



\begin{frame}
\frametitle{Driving mechanisms and further questions} 
What explains (or does not explain) these changes?
\begin{itemize}\setlength\itemsep{3pt}
\item Possible confounders from municipality-wide forces? No effects on private schools
\item No significant declines in educational spending in turnover municipalities
\item Not related to political ideologies: Declines observed in both left/right wing parties
\item Experience and capability of new personnel $\downarrow\to$ trade-off productivity and loyalty
  \end{itemize}\medskip
What remains unanswered?
\begin{itemize}\setlength\itemsep{3pt}
\item Politically motivated replacements can be a costly byproduct on public services
\begin{itemize}\setlength\itemsep{3pt}
\item Even in closely-contested, fair elections
\end{itemize}
\item However, discretion may be needed to weed out low-performing, entrenched workers
\item So what determines the relative strengths of these opposing forces?
\end{itemize}
\end{frame}



\subsection{Fenizia (2022): Managers and Productivity in Public Sector}
\begin{transitionframe}
\begin{center}
  \begin{huge}
    \textcolor{NTHU1}{%
Fenizia (2022):\\ Managers and Productivity in Public Sector
    }
  \end{huge}
\end{center}
\end{transitionframe}



\begin{frame}
\frametitle{Q: Do talented managers in public sector improve performance?} 
How important are managers in public sectors?
\begin{itemize}\setlength\itemsep{3pt}
\item Oversee day-to-day operations and supervise policy implementation 
\item Limited tools available compared to private sector managers`'
\begin{itemize}
\item Limited discretion in firing and promoting workers
\end{itemize}
\item How to measure performance of government agencies?
\end{itemize}\medskip
This paper uses innovative measures of performance to study the role of managers
\begin{itemize}\setlength\itemsep{3pt}
\item Italian Social Security Agency data (INPS): Covering varous welfare programs
\item Uses welfare claims as measure of output (homogenous across different offices)
\item All sites have same rules and physical capital$\to$ distinguish manager contributions
\end{itemize}\medskip
\end{frame}

\begin{frame}
\frametitle{Background and data} 
Role of INPS 
\begin{itemize}\setlength\itemsep{3pt}
\item Administers all social welfare and insurance programs in Italy 
\item Claims are processed in local offices run by single managers 
\item Managers allocate work, determine workplace practices, but cannot fire workers
\item Managers selected among limited pool of candidates and rotate every 5 years 
\end{itemize}\medskip
Data sources on managers (2011Q1 - 2017Q2)
\begin{itemize}\setlength\itemsep{3pt}
\item Output: $\sum$ claims weighed by complexity (procesing time) per worker
\item Service quality: Timeliness and error rate
\item Worker characteristics: Job title, age, office locations etc
\end{itemize}
\end{frame}

\begin{frame}
\frametitle{Empirical strategies} 
Do productive managers improve performance?
\begin{itemize}\setlength\itemsep{3pt}
\item Fixed effects regression with office($i$), time($t$), and manager($m(i,t)$) fixed effects
\[
\log{P}_{it} = \alpha_i + \tau_t + \theta_{m(i,t)}+u_{it}
\]\vspace{-20pt}
\item Requires that manager moves are exogenous conditional on office and time effects
\begin{itemize}
  \item Observable characteristics should not predict manager quality
\item Violated if better managers are assigned to worse offices
\end{itemize}
\item Compare adjusted $R^2$ to test how much of output variation is explained by managers 
\end{itemize}\medskip
What makes for a productive manager?
\begin{itemize}\setlength\itemsep{3pt}
\item Regress changes in outcomes with quarter before manager change to talent changes
\[
y_{i,t+k}-y_{i,t-1} = \pi_0^k + \pi_1^k(\hat{\theta}_{i,new}^k-\hat{\theta}_{i,old})+\gamma^k X_{i}+ (e_{i,t+k}-e_{i,t-1})
\]\vspace{-20pt}

\end{itemize}
\end{frame}

\begin{frame}
\frametitle{Managers explain 10\% of variation in productivity } 
\begin{figure}
  \includegraphics[keepaspectratio, width=0.57\textwidth]{f2022_t1.png}
\end{figure}
\begin{itemize}
\item Adjusted $R^2$: 0.68 in (2) $\to$ 0.76 in (3): 8 p.p increase
\item $0.08/0.762\simeq 0.104$: Manager explains 10\% of variation in productivity
\item Robust to selective sorting, other ways of variance decomposition
\end{itemize}

\end{frame}


\begin{frame}
\frametitle{Talented managers improve productivity, with workers decreasing...} 
\begin{figure}
  \includegraphics[keepaspectratio,width=0.6\textwidth]{f2022_f1.png}
\end{figure}
\end{frame}

\begin{frame}
\frametitle{Talented managers induced by retirement of older workers} 
\begin{figure}
  \includegraphics[keepaspectratio,width=0.6\textwidth]{f2022_t2.png}
\end{figure}
\end{frame}

\begin{frame}
\frametitle{Takeaways and discussions}
Other outcomes 
\begin{itemize}\setlength\itemsep{3pt}
\item Trade-off in quality? Not observed
\item Outcomes not driven by shifting towards easier claims to handle
\item Thus managers do have a meaningful impact on productivity
\end{itemize}\medskip
Remaining questions
\begin{itemize}\setlength\itemsep{3pt}
\item The managers oversee relatively small teams with well-defined goals this context?
\begin{itemize}\setlength\itemsep{3pt}
\item What about large projects and teams?
\item What about organization with no fixed roles and less structure?
\end{itemize}
\item Other factors may determine managers' role in productivity
\begin{itemize}\setlength\itemsep{3pt}
\item Ability to obtain inputs efficiently (no bribing, better communication)
\item Ability to attract more talented workers under the command 
\end{itemize}
\end{itemize}\medskip
\end{frame}


\section{Recruiting capable bureaucrats to public sectors}
\subsection{Dal Bó, Finan, Rossi (2013): Strengthening State Capablities: The Role of Financial Incentives in the Call to Public Service}
\begin{transitionframe}
\begin{center}
  \begin{huge}
    \textcolor{NTHU1}{%
Dal Bó, Finan, Rossi (2013):\\ Strengthening State Capablities: The Role of Financial Incentives in the Call to Public Service
    }
  \end{huge}
\end{center}
\end{transitionframe}


\begin{frame}
\frametitle{Q: Can fiscal incentives attract capable individuals to public sector?} 
Professionalized bureaucracy require individuals with capabilities and integrity
\begin{itemize}\setlength\itemsep{3pt}
\item These do not always co-move
\begin{itemize}\setlength\itemsep{3pt}
\item Capable individuals may seek opportunities in private sectror (higher pay)
\item This may leave public sector with less able individuals
\item Not all capable individuals have integrity:
\end{itemize}
\item So what makes up motivated agents and how to attract them to public sector?
\end{itemize}\medskip
This paper investigates role of incentives while investigating various aspects of quality
\begin{itemize}\setlength\itemsep{3pt}
\item How do monetary and non-monetary incentives attract candidates?
\item Estimates if higher wages attract more candidates and those of higher quality
\item Also estimates the labor supply elasticity facing public sector in Mexico
\item Heterogeneity by locational characteristics (local crime rates)
\end{itemize}
\end{frame}



\begin{frame}
\frametitle{Context of the experiment}
Mexico's regional development program (RDP) in 2011
\begin{itemize}\setlength\itemsep{3pt}
  \item Network of 350 agents \& 50 coordinators to identify needs, report to federal govt
\item Exogenously assign different wage offers across different posts (5000 vs 3750 MXP)
\item Candidates screened to document various characteristics
\begin{itemize}\setlength\itemsep{3pt}
  \item Aptitude: intelligence  (IQ) and other qualities known to determine market value
  \item Motivation: Integrity, prosocial inclinations, public service motivation (interviews)
  \item Attraction \& commitment to policy making, justice, civicness, compassion, self-sacrifice
\end{itemize}
\item Vacancies randomized to one of 4 strata $\to$ Offers randomized to indiv. in that stratum
\begin{itemize}\setlength\itemsep{3pt}
  \item High vs low wage, High ($>9$) vs normal IQ (7-9 Ravens matrix)
\end{itemize}
\end{itemize}\medskip
Estimating causal effects: OLS (due to random assignment design)
\begin{itemize}\setlength\itemsep{3pt}
  \item Effect on character $Y$ in applicant pool (region): $Y_{icr}=\beta_1 T_c + \xi_r + e_{icr}$
  \item Effect on acceptance $A$ among applicants (strata): $A_{ics}=\gamma_1 T_c + \beta X_i + \xi_s + e_{ics}$
  \item Hetergeneity by municipalities: $A_{ics}=\gamma_1 T_c + \gamma_2 T_c \times W_m + \gamma_3 W_m + \beta X_i + \xi_s + e_{ics}$
\end{itemize}
\end{frame}

\begin{frame}
\frametitle{Higher wage $\to$ High ability applicants $\Uparrow$} 
\begin{columns}
\begin{column}{0.5\textwidth}
\begin{figure}
  \includegraphics[keepaspectratio, width=\textwidth]{dfr2013_t1_1.png}
\end{figure}
\end{column}
\begin{column}{0.5\textwidth}
\begin{figure}
  \includegraphics[keepaspectratio, width=\textwidth]{dfr2013_t1_2.png}
\end{figure}
\end{column}
\end{columns}
\end{frame}

\begin{frame}
\frametitle{Higher wage $\to$ Better motivated applicants $\Uparrow$} 
\begin{columns}
\begin{column}{0.5\textwidth}
\begin{figure}
  \includegraphics[keepaspectratio, width=\textwidth]{dfr2013_t2_1.png}
\end{figure}
\end{column}
\begin{column}{0.5\textwidth}
\begin{figure}
  \includegraphics[keepaspectratio, width=\textwidth]{dfr2013_t2_2.png}
\end{figure}
\end{column}
\end{columns}
\end{frame}

\begin{frame}
\frametitle{Higher wage $\to$ Applicants more likely to accept} 
\begin{figure}
  \includegraphics[keepaspectratio, width=0.7\textwidth]{dfr2013_t3.png}
\end{figure}
\end{frame}

\begin{frame}
\frametitle{Higher wage needed to offset disadvantageous regional traits} 
\begin{figure}
  \includegraphics[keepaspectratio, width=0.7\textwidth]{dfr2013_t4.png}
\end{figure}
\end{frame}

\begin{frame}
\frametitle{Additional findings and takeaways}
Labor supply elasticity $(\epsilon)$ for public sector with respect to wages 
\begin{itemize}\setlength\itemsep{3pt}
\item $\epsilon  = \epsilon_{\text{pool}}  + \epsilon_{\text{accept}}+\epsilon_{\text{pool}} \times \epsilon_{\text{accept}}\times\text{wage difference}$
\item Wage difference: 33\% (=(5000-3750)/3750)
\item Pool: 4.7 person $\uparrow$ from 18.1 person (26.3\%$\uparrow$) (Table 3) $\to$ Elasticity: 0.8 (=26.3/33)
\item Acceptance:  15.1 p.p $\uparrow$ from 42.9\% (35.2\% $\uparrow$) (Table 5) $\to$ Elasticity: 1.07 (=35.2/33)
\item $0.8+1.07+0.8\times1.07\times0.33=2.15$ ( 1\% $\uparrow$ in wages $\to$ Supply  2.15\% $\uparrow$)
\end{itemize}\medskip
Takeaways
\begin{itemize}\setlength\itemsep{3pt}
\item Offering high wages attracts better and motivated candidates
\item However, findings hold only if motivation and ability are positively correlated
 \begin{itemize}
\item What if service commitment is sacrificed for higher pay (prioritize superiors over citizens?)
\end{itemize}
\item Silent on post-acceptance performance: High pay to keep people motivated? 
\end{itemize}
\end{frame}



\subsection{Desserano (2019): Financial Incentives as Signals: Experimental Evidence from the Recruitment of Village Promoters in Uganda}
\begin{transitionframe}
\begin{center}
  \begin{huge}
    \textcolor{NTHU1}{%
Desserano (2019):\\ Financial Incentives as Signals: Experimental Evidence from the Recruitment of Village Promoters in Uganda
    }
  \end{huge}
\end{center}
\end{transitionframe}



\begin{frame}
\frametitle{Q: How does financial incentives shape perception on public posting?} 
What does financial incentives signal to applicants when informaion is incomplete?
\begin{itemize}\setlength\itemsep{3pt}
 \item Information about the job is incomplete in many instances
\item Financial incentives may provide signal on the nature of the task and the organization
 \begin{itemize}\setlength\itemsep{3pt}
\item High pay $\to$ More demanding and less enjoyable
\item May signal exploitative intentions
\end{itemize}
\item How does this apply to public positions that require pro-social characteristics?
\end{itemize}\medskip
This paper tests the effect of fiscal incentives on job perceptions and application decisions
\begin{itemize}\setlength\itemsep{3pt}
  \item Experiment with new position involving social and business task in Uganda
  \item Financial incentives signal the business-oriented nature
  \item Affects composition and size of the appicant pool
   \begin{itemize}\setlength\itemsep{3pt}
\item More applicants, but discourage applicants with pro-social preferences
\end{itemize}
\item Thus, financial incentives may crowd out pro-social motives under incomplete info
  \end{itemize}
\end{frame}



\begin{frame}
\frametitle{Experimental setting}
Recruitment of community health promoters (CHP) in Uganda
\begin{itemize}\setlength\itemsep{3pt}
    \item Hire unskilled local agents and teach them how to serve their own villages
    \item Social component: Promote health education
    \item Business side: Sell household products at low prices
    \item Completely new: No prior on relative importance, difficulty, recruiter's goals etc
\end{itemize}\medskip
The experiment creates variations in expected earnings (fix actual incentives)
\begin{itemize}\setlength\itemsep{3pt}
  \item Most CHPs talk together: Unequal actual incentives could cause problems
  \item Observe how differences in expectations affect types of applicants at selection
  \item How this was implemented
  \begin{itemize}\setlength\itemsep{3pt}
  \item Applicants are told that earnings vary across workers (profit margin for selling)
  \item Randomly given a different point of the income distribution (max, avg, min)  
  \item This defines high/medium/low-pay treatments
  \end{itemize}
\end{itemize}

\end{frame}



\begin{frame}
  \frametitle{Experiment details}
  \begin{figure}
    \centering
    \includegraphics[keepaspectratio, width=0.7\textwidth]{d2019_f1.png}
  \end{figure}
\end{frame}



\begin{frame}
\frametitle{Empirical strategy} 
Recruitment and information experiments
\begin{itemize}\setlength\itemsep{3pt}
  \item Recruitment: Advertised to 4,863 eligible candidates in 315 communities
\begin{itemize}\setlength\itemsep{3pt}
  \item Treatment adminstered to microfinance groups and interviewed on intent to apply 
\end{itemize}
  \item Information: 6,844 respondents asked to report their perception of the job
\begin{itemize}\setlength\itemsep{3pt}
  \item Gathered socioeconomic questions and pro-social preference information
  \item 1/2 asked on their expected earnings, 1/2 asked on their perception
\end{itemize}
\item Prosociality: Volunteer experience, and job feature the applicant prioritizes
\end{itemize}\medskip
Regression analysis
\begin{itemize}\setlength\itemsep{3pt}
    \item Information treatment
    \[
    Y_i = \alpha + \beta_1\text{Medium}_i+\beta_2\text{High}_i+X_i\eta + e_i
    \]\vspace{-20pt}
    \item Recruitment treatment: Interact Medium and High with pro-social traits
    \begin{itemize}\setlength\itemsep{3pt}
    \item One version with microfinance group FE
    \item Other with NGO branch FE (with dummy for trait, non-trait)
\end{itemize}
  \end{itemize}
\end{frame}



\begin{frame}
  \frametitle{Financial incentives signal priority on private goals (vs. communal)}
  \begin{figure}
    \centering
    \includegraphics[keepaspectratio, width=1\textwidth]{d2019_t1.png}
  \end{figure}
\end{frame}

\begin{frame}
  \frametitle{Prosociality declines in high-pay treatment}
  \begin{figure}
    \centering
    \includegraphics[keepaspectratio, width=1\textwidth]{d2019_t2.png}
  \end{figure}
\end{frame}



\begin{frame}
\frametitle{Other findings and remaining questions} 
Other findings: Financial incentives may crowd-out pro-social preferences
\begin{itemize}\setlength\itemsep{3pt}
  \item The size of the pool is largest in high-pay treatment
  \item However, those recruited under high-pay remain less pro-social even after working
\item Higher retention of CHP workers in low-pay treatment
\end{itemize}\medskip
Remaining questions
\begin{itemize}\setlength\itemsep{3pt}
    \item Financial incentives highlights business aspects of the work, even after recruitment
    \begin{itemize}\setlength\itemsep{3pt}
      \item Any way to enhance pro-social motives afterwards?
    \end{itemize}
    \item Different setting compared to Dal Bó, Finan, Rossi (2013)
     \begin{itemize}\setlength\itemsep{3pt}
      \item Known job of public workers vs uncertain characteristics of CHPs
      \item Financial incentives may not have a signalling effect in the former
      \item Thus, the crowding-out effect may not be observed in some cases
    \end{itemize}
\end{itemize}
  
\end{frame}


\subsection{Moreira, Pérez (2024): Civil Service Exams and Organizational Performance: Evidence from the Pendleton Act}
\begin{transitionframe}
\begin{center}
  \begin{huge}
    \textcolor{NTHU1}{%
Moreira, Pérez (2024):\\ Civil Service Exams and Organizational Performance: Evidence from the Pendleton Act
    }
  \end{huge}
\end{center}
\end{transitionframe}

\begin{frame}
\frametitle{Q: How do civil service exams affect organizational performance?}
Professionalized bureaucracy increasingly seen as key for economic development
\begin{itemize}\setlength\itemsep{3pt}
\item Many tries to achieve this by Introducing competitive exams
\item The intention is to limit favoritism and help select more qualified empoyees
\item But additional positions not covered by exams could distort hierarchical structure
\item So the effect on performance may be ambiguous event with better selection
\end{itemize}\medskip
This paper studies impact of civil service exams on organizational performance in the US
\begin{itemize}\setlength\itemsep{3pt}
\item Using 1883 Pendleton Civil Service Reform Act (mandated exams for federal workers)
\item Focus on Customs Service (responsible for large share of federal revenues)
\item Despite improved selection and retention, organizational performance did not improve
\item Increases in exam-exempted positions and retentions of managers explains findings
\end{itemize}
\end{frame}


\begin{frame}
\frametitle{Historical Background}
Customs revenue collection in historic US
\begin{itemize}\setlength\itemsep{3pt}
\item Customs revenue accounted for more than half the federal revenues in 1880s
\item Hiring was based on political and personal connections prior to reforms
\item Inefficient practies: Cost of collecting \$1 larger in US than DEU, ENG, FRA 
\begin{itemize}\setlength\itemsep{3pt}
\item Largely due to overstaffing and `errors and frauds' of these employees
\end{itemize}
\end{itemize}\medskip
Implementation of the 1883 Pendleton reforms on customs
\begin{itemize}\setlength\itemsep{3pt}
\item Applied to districts with 50+ employees and positions earning \$900+ by 1883
  \begin{itemize}\setlength\itemsep{3pt}
\item These districts are called `classified' positions
\item Only allowed to select candidates among top 4 scores in open exams
\end{itemize}
\item Tested on practical knowledge relevant to the position rather than academic training
\item Other components applied to ALL districts (making treated vs control comparable)
  \begin{itemize}\setlength\itemsep{3pt}
\item No political assessments and hiring individuals `using intoxicating bevarages to excess'
\end{itemize}
\end{itemize}\medskip
\end{frame}





\begin{frame}
  \frametitle{Data and empirical strategy}
Data sources: Covers 1871-1893 (Authors digitized the records themselves)
\begin{itemize}\setlength\itemsep{3pt}
\item Personnel data: Official Registers of the US (biennial) linked to US population census
\begin{itemize}\setlength\itemsep{3pt}
\item Name, job title, state/country of birth, compensation, place of employment
\end{itemize}
\item Finance records: Annual Report of Secretary of the Treasury on State of the Finances
\begin{itemize}\setlength\itemsep{3pt}
\item Revenues and expenditures of different segments of government agencies
\end{itemize}
\end{itemize}\medskip
Leverage difference-in-differences induced by 1883 reforms 
\[
Y_{ct}=\alpha_c +\alpha_t + \beta\text{Classified}_c\times\text{After}_t+\gamma X_{ct}+e_{ct}
\]\vspace{-20pt}
\begin{itemize}\setlength\itemsep{3pt}
\item In some regressions, $\beta$ is broken down across different time periods
\item Authors verify that outcomes evolve in parallel pattern pre-reform (parallel trends)
\item No evidence of districts manipulating their size downward to avoid reforms
\end{itemize}
\end{frame}


\begin{frame}
\frametitle{Exams $\to$ Hiring of those with professional backgrounds} 
\begin{figure}
  \includegraphics[keepaspectratio, width=0.7\textwidth]{mp2024_t1.png}
\end{figure}
\end{frame}

\begin{frame}
\frametitle{Less likely to leave after reform} 
\begin{figure}
  \includegraphics[keepaspectratio, width=0.7\textwidth]{mp2024_t2.png}
\end{figure}
\end{frame}

\begin{frame}
\frametitle{However, reforms did not lead to better cost-effectiveness} 
\begin{figure}
  \includegraphics[keepaspectratio, width=0.67\textwidth]{mp2024_f1.png}
\end{figure}
\end{frame}


\begin{frame}
\frametitle{Mechanisms and unanswered questions}
The incomplete reach of the reforms
\begin{itemize}\setlength\itemsep{3pt}
\item Managers and positions with salaries below a threshold not subject to reforms
\item The number of low-paid positions inflated following reforms (twice)
\item Selection of managers were left unchanged 
\end{itemize}\medskip

Room for further research on other aspects
\begin{itemize}\setlength\itemsep{3pt}
\item Reducing favoritism in hiring can increase retention and quality in selection
\item But does not automatically lead to better practices
\begin{itemize}\setlength\itemsep{3pt}
\item Need to prevent responses that can offset the effects of reforms
\item Managers also matter - Need to minimize incomplete reach
\end{itemize}
\item How to design organizational reform that minimizes these side effects?
\end{itemize}
\end{frame}


\begin{frame}
\frametitle{Remaining questions for future research}
Attracting capable and intrinsically-motivated individuals 
\begin{itemize}\setlength\itemsep{3pt}
\item Higher wages attract capability but may crowd out pro-social preferences
\item Non-monetary incentives matter but hard to design at scale
\end{itemize}\medskip
 When does discretion improve selection vs. open the door to patronage?
\begin{itemize}\setlength\itemsep{3pt}
\item Optimal mix of rule-based and discretionary selection likely varies by context
\item How do political institutions shape this tradeoff?
\end{itemize}\medskip

 Does better selection translate into better performance?
\begin{itemize}\setlength\itemsep{3pt}
\item Selection is only one part: Need to ensure employees are incentivized afterwrads
\item How to design organizational structure that maximizes incentives post-selection?
\end{itemize}
\end{frame}



%%%%%%%%%%%
\end{document}
